% Môn: Vật lý
% Lớp: 11
% Chương: Chương IV. Dòng điện. Mạch điện
% Bài: L11C4  Bài 24. Nguồn điện
% Số câu: 25
% App đọc file này trực tiếp — sửa rồi Commit trên GitHub, bấm Đồng bộ GitHub .tex

% L11C4  Bài 24. Nguồn điện

% ===== Câu 1 =====
% ID: L11C4B24-01-TL
% Mức: NB
\dangbt{Chưa có dạng}
\begin{ex}
% ID: L11C4B24-01-TL
% Mức: NB
Suất điện động của một nguồn điện là $1,5\,\mathrm{V}$, lực lạ làm di chuyển điện tích thực hiện một công 6 mJ. Lượng điện tích dịch chuyển khi đó là A. 15 $\cdot$ 10^{-3} C. B. 4 $\cdot$  10^{-3} C. C. 0,5 $\cdot$  10^{-3} C. D. 1,5 $\cdot$  10^{-3} C. $\nguon{ly11 kntt.pdf}$
\choiceTF

\loigiai{
Chọn B.

Lượng điện tích dịch chuyển $q = \dfrac{A}{\text{E}} = \dfrac{6\cdot 10^{-3}}{1,5} = 4\cdot 10^{-3}C$
}
\end{ex}

% ===== Câu 2 =====
% ID: L11C4B24-01-TN
% Mức: NB
\begin{ex}
% ID: L11C4B24-01-TN
% Mức: NB
Kết luận nào sau đây đúng khi nói về tác dụng của nguồn điện? $\nguon{ly11 kntt.pdf}$
\choice
{\True dùng để tạo ra và duy trì hiệu điện thế nhằm duy trì dòng điện trong mạch.}
{dùng để tạo ra các ion âm.}
{dùng để tạo ra các ion dương.}
{dùng để tạo ra các ion âm chạy trong vật dẫn.}
\loigiai{
Chọn A.

Nguồn điện dùng để tạo ra và duy trì hiệu điện thế nhằm duy trì dòng điện trong mạch.
}
\end{ex}

% ===== Câu 3 =====
% ID: L11C4B24-02-TL
% Mức: TH
\begin{bt}
% ID: L11C4B24-02-TL
% Mức: TH
Suất điện động của một nguồn điện là $12\,\mathrm{V}$. Tính công của lực lạ khi dịch chuyển một lượng điện tich là $0,5\,\mathrm{C}$ bên trong nguồn điện từ cực âm đến cực dương của nó. $\nguon{ly11 kntt.pdf}$
\loigiai{
Công của lực lạ: $A= qE = 0,5.12 =6\,\mathrm{J}$.
}
\end{bt}

% ===== Câu 4 =====
% ID: L11C4B24-02-TN
% Mức: NB
\begin{ex}
% ID: L11C4B24-02-TN
% Mức: NB
Kết luận nào sau đây sai khi nói về suất điện động của nguồn điện? $\nguon{ly11 kntt.pdf}$
\choice
{Suất điện động của nguồn điện đặc trưng cho khả năng thực hiện công của nguồn điện. $A$}
{Suất điện động của nguồn điện được đo bằng thương số q.}
{Đơn vị của suất điện động là vôn (V).}
{\True Suất điện động của nguồn điện đặc trưng cho khả năng tích điện của nguồn điện.}
\loigiai{
Chọn D.

$D$ – sai vì suất điện động của nguồn điện đặc trưng cho khả năng thực hiện công của nguồn điện.
}
\end{ex}

% ===== Câu 5 =====
% ID: L11C4B24-03-TL
% Mức: TH
\begin{ex}
% ID: L11C4B24-03-TL
% Mức: TH
Một acquy có suất điện động $6\,\mathrm{V}$, sản ra một công là $360\,\mathrm{J}$ khi acquy này phát điện trong 5 phút. a) Tính lượng điện tích dịch chuyển trong acquy. b) Tính cường độ dòng điện chạy qua acquy. $\nguon{ly11 kntt.pdf}$
\choiceTF

\loigiai{
a) Điện lượng dịch chuyển trong acquy là: $. \text{E} = \dfrac{A}{q}\Rightarrow q = \dfrac{A}{\text{E}} = \dfrac{360}{6} =$ 60 $\mathrm{C}$ \right. $.

b) Cường độ dòng điện chạy qua acquy là:$ I = \dfrac{q}{\text{\Delta }t} = \dfrac{60}{5.60} = 0,2\text{A}.
}
\end{ex}

% ===== Câu 6 =====
% ID: L11C4B24-03-TN
% Mức: NB
\begin{ex}
% ID: L11C4B24-03-TN
% Mức: NB
Khi nói về nguồn điện, phát biểu nào dưới đây sai? $\nguon{ly11 kntt.pdf}$
\choice
{Mỗi nguồn có hai cực luôn ở trạng thái nhiễm điện khác nhau.}
{Nguồn điện là cơ cấu để tạo ra và duy trì hiệu điện thế nhằm duy trì dòng điện trong đoạn mạch.}
{Để tạo ra các cực nhiễm điện, cần phải có lực thực hiện công tách và chuyển các electron hoặc ion dương ra khỏi điện cực, lực này gọi là lực lạ.}
{\True Nguồn điện là pin có lực lạ là lực tĩnh điện.}
\loigiai{
Chọn D.

$D$ – sai vì lực lạ không phải là lực tính điện.
}
\end{ex}

% ===== Câu 7 =====
% ID: L11C4B24-04-TL
% Mức: TH
\begin{ex}
% ID: L11C4B24-04-TL
% Mức: TH
Một bộ acquy đầy điện có thể cung cấp dòng điện $4\,\mathrm{A}$ liên tục trong 2 giờ thì phải nạp lại. a) Tính cường độ dòng điện mà acquy này có thể cung cấp liên tục trong 40 giờ thì phải nạp lại. b) Tính suất điện động của acquy nếu trong thời gian hoạt động trên đây, nó $\sin$ h ra một công là 172,8 kJ. BÀ̀ $I$ 25. NĂNG LƯỢNG VÀ CÔNG SUẤ $T$ ĐIỆ $N\nguon{ly11 kntt.pdf}$
\choiceTF

\loigiai{
a) Điện lượng: q = It = 4.2.3600 = $28800\,\mathrm{C}$

Cường độ dòng điện mà acquy này có thể cung cấp liên tục trong 40 giờ: $I^{'} = \dfrac{q}{t^{'}} = \dfrac{28800}{40.3600} = 0,$ 2 $\mathrm{A}$ \text{.} $b) Tính suất điện động của acquy:$ \text{E} = \dfrac{A}{q} = \dfrac{172,8.1000}{28800} = $6$ \mathrm{V}.
}
\end{ex}

% ===== Câu 8 =====
% ID: L11C4B24-04-TN
% Mức: NB
\begin{ex}
% ID: L11C4B24-04-TN
% Mức: NB
Nguồn điện tạo ra hiệu điện thế giữa hai cực bằng cách $\nguon{ly11 kntt.pdf}$
\choice
{\True tách electron ra khỏi nguyên tử và chuyển electron và ion ra khỏi các cực của nguồn.}
{$\sin$ h ra ion dương ở cực âm.}
{$\sin$ h ra electron ở cực dương.}
{làm biến mất electron ở cực dương.}
\loigiai{
Chọn A.

Nguồn điện tạo ra hiệu điện thế giữa hai cực bằng cách tách electron ra khỏi nguyên tử và chuyển electron và ion ra khỏi các cực của nguồn.
}
\end{ex}

% ===== Câu 9 =====
% ID: L11C4B24-05-TL
% Mức: TH
\begin{bt}
% ID: L11C4B24-05-TL
% Mức: TH
Một nguồn điện có điện trở trong $r=0,1\,\Omega$ được mắc nối tiếp với điện trở $R=4,8\,\Omega$ tạo thành mạch kín. Hiệu điện thế giữa hai cực của nguồn là $U=12\,\mathrm{V}$. Tính suất điện động của nguồn và cường độ dòng điện chạy trong mạch. $\nguon{ly11 kntt.pdf}$
\loigiai{
Vì hiệu điện thế mạch ngoài là $U=IR$ nên
$I=\dfrac{U}{R}=\dfrac{12}{4,8}=2,5 \mathrm{A}.$
Suất điện động của nguồn là
$\mathcal{E}=U+Ir=12+2,5\cdot0,1=12,25 \mathrm{V}.$
}
\end{bt}

% ===== Câu 10 =====
% ID: L11C4B24-05-TN
% Mức: NB
\begin{ex}
% ID: L11C4B24-05-TN
% Mức: NB
Câu nào sau đây sai? $\nguon{ly11 kntt.pdf}$
\choice
{Suất điện động của nguồn điện là đại lượng đặc trưng cho khả năng $\sin$ h công của nguồn điện.}
{\True Suất điện động của nguồn điện được xác định bằng công suât dịch chuyển vòng kín của mạch điện.}
{Suất điện động của nguồn điện bằng công để di chuyển điện tích dương $1\,\mathrm{C}$ từ cực âm đến cực dương bên trong nguồn.}
{Suất điện động được đo bằng thương số giữa công $A$ của lực lạ để di chuyển một điện tích dương q từ cực âm đến cực dương bên trong nguồn điện và độ lớn của điện tích đó.}
\loigiai{
Chọn B.

$B$ - sai
}
\end{ex}

% ===== Câu 11 =====
% ID: L11C4B24-06-TL
% Mức: TH
\begin{bt}
% ID: L11C4B24-06-TL
% Mức: TH
Mạch kín gồm nguồn điện có suất điện động $\mathcal{E} = 24\$, $\mathrm{V},$ điện trở trong $r = 0$,5 $\Omega$ và hai điện trở $R_1 = 10\Omega$; $R_2 = 50\Omega$ mắc nối tiếp. Một vôn kế mắc song song với điện trở R_2, chỉ giá trị $16\,\mathrm{V}$. Tìm điện trở của vôn kế. $\nguon{ly11 kntt.pdf}$
\loigiai{
Điện trở mạch ngoài: $R = \dfrac{R_{2}R_{V}}{R_{2} + R_{V}} + R_{1}$.

Điện trở toàn mạch bằng $R + r = \dfrac{R_{2}R_{V}}{R_{2} + R_{V}} + R_{1} + r$.

Cường độ dòng điện trong mạch $I = \dfrac{\text{E}}{R + r}$.

Số chỉ vôn kế $U_{V} = I \cdot \dfrac{R_{2}R_{V}}{R_{2} + R_{V}}$.

Vậy ta có: $16 = \dfrac{\text{E}}{R_{1} + \dfrac{R_{2}R_{V}}{R_{2} + R_{V}} + r} \cdot \dfrac{R_{2}R_{V}}{R_{2} + R_{V}}$.

Thay số $\Rightarrow$ R_(V)=13,3 $\Omega$ .
}
\end{bt}

% ===== Câu 12 =====
% ID: L11C4B24-06-TN
% Mức: NB
\begin{ex}
% ID: L11C4B24-06-TN
% Mức: NB
Công của nguồn điện là $\nguon{ly11 kntt.pdf}$
\choice
{lượng điện tích mà nguồn điện $\sin$ h ra trong $1\,\mathrm{s}$.}
{\True công của lực lạ làm dịch chuyển điện tích bên trong nguồn.}
{công của dòng điện trong mạch kín $\sin$ h ra trong $1\,\mathrm{s}$.}
{công của dòng điện khi dịch chuyển một đơn vị điện tích trong mạch kín.}
\loigiai{
Chọn B.

Công của nguồn điện là công của lực lạ làm dịch chuyển điện tích bên trong nguồn.
}
\end{ex}

% ===== Câu 13 =====
% ID: L11C4B24-07-TL
% Mức: TH
\begin{bt}
% ID: L11C4B24-07-TL
% Mức: TH
Hai điện trở $R_1 = 2\Omega$, $R$ 2 = 6 $\Omega$ mắc vào nguồn điện có suất điện động $\mathcal{E}$ và điện trở trong r. Khi R_1, R_2 mắc nối tiếp, cường độ dòng điện trong mạch chính $I = 0$, $5\,\mathrm{A}$.Khi R_1, $R$ 2 mắc song song, cường độ dòng điện trong mạch chính $I = 1$, $8\,\mathrm{A}$.Tìm giá trị của suất điện động $\mathcal{E}$ và điện trở trong r. $\nguon{ly11 kntt.pdf}$
\loigiai{
- Khi R₁ mắc nối tiếp với R₂, điện trở toàn mạch là: $R= R₁ + R₂ + r = 8 + r.\left. \Rightarrow I = \dfrac{\text{E}}{R} = \dfrac{\text{E}}{8 + r}\Rightarrow 0,5 = \dfrac{\text{E}}{8 + r} \right.$ (1)

Khi R₁//R₂, điện trở toàn mạch là: $R^{'} = \dfrac{R_{1}R_{2}}{R_{1} + R_{2}} + r = \dfrac{12}{8} + r = 1,5 + r\left. \Rightarrow I^{'} = \dfrac{\text{E}}{R^{'}} = \dfrac{\text{E}}{1,5 + r}\Rightarrow 1,8 = \dfrac{\text{E}}{1,5 + r} \right.$ (2)

Từ (1) và (2) ta suy ra: r = 1 $\Omega$ ; $E$ = 4, $5\,\mathrm{V}$.
}
\end{bt}

% ===== Câu 14 =====
% ID: L11C4B24-07-TN
% Mức: NB
\begin{ex}
% ID: L11C4B24-07-TN
% Mức: NB
Suất điện động của nguồn điện một chiều là $\mathcal{E} = 4\,\mathrm{V}$. Công của lực lạ làm dịch chuyển một lượng điện tích $q = 5\,\mathrm{mC}$ giữa hai cực bên trong nguồn điện là
\choice
{1,5 mJ.}
{0,8 mJ.}
{\True 20 mJ.}
{5 mJ.}
\loigiai{
Chọn C.

Công cần tìm $A= qE = 5\cdot10^{-3}.4 = 0,02\,\mathrm{J}$ = 20mJ
}
\end{ex}

% ===== Câu 15 =====
% ID: L11C4B24-08-TN
% Mức: NB
\begin{ex}
% ID: L11C4B24-08-TN
% Mức: NB
Một acquy có suất điện động là $12\,\mathrm{V}$, $\sin$ h ra công là $720\,\mathrm{J}$ đễ duy trì dòng điện trong mạch trong thời gian 1 phút. Cường độ dòng điện chạy qua acquy khi đó là $\nguon{ly11 kntt.pdf}$
\choice
{$I = 1$, $2\,\mathrm{A}$.}
{$I = 5$, $0\,\mathrm{A}$.}
{\True 1 = $0,2\,\mathrm{A}$.}
{$I = 2$, $4\,\mathrm{A}$.}
\loigiai{
Chọn C.

Cường độ dòng điện $I = \dfrac{A}{Et} = \dfrac{720}{12.60} =$ 1 $\mathrm{A}$
}
\end{ex}

% ===== Câu 16 =====
% ID: L11C4B24-09-TN
% Mức: NB
\begin{ex}
% ID: L11C4B24-09-TN
% Mức: NB
Một acquy đầy điện có dung lượng $20\,\mathrm{A}$.h. Biết cường độ dòng điện mà nó cung cấp là $0,5\,\mathrm{A}$. Thời gian sử dụng của acquy là $\nguon{ly11 kntt.pdf}$
\choice
{$t = 5\$, $\mathrm{h}$.}
{\True $t = 40\$, $\mathrm{h}$.}
{$t = 20\$, $\mathrm{h}$.}
{$t = 50\$, $\mathrm{h}$.}
\loigiai{
Chọn B.

Thời gian cần tìm $t = \dfrac{20}{0,5} = 40h$
}
\end{ex}

% ===== Câu 17 =====
% ID: L11C4B24-10-TN
% Mức: NB
\begin{ex}
% ID: L11C4B24-10-TN
% Mức: NB
Suất điện động của nguồn điện là đại lượng được đo bằng $\nguon{ly11 kntt.pdf}$
\choice
{công của lực lạ tác dụng lên điện tích q dương.}
{thương số giữa công và lực lạ tác dụng lên điện tích q dương.}
{thương số giữa lực lạ tác dụng lên điện tích q dương và độ lớn điện tích ấy.}
{\True thương số giữa công của lực lạ dịch chuyển điện tích dương q từ cực âm đến cực dương trong nguồn và độ lớn của điện tích đó.}
\loigiai{
Chọn D.

Suất điện động của nguồn điện là đại lượng được đo bằng thương số giữa công của lực lạ dịch chuyển điện tích dương q từ cực âm đến cực dương trong nguồn và độ lớn của điện tích đó
}
\end{ex}

% ===== Câu 18 =====
% ID: L11C4B24-11-TN
% Mức: NB
\begin{ex}
% ID: L11C4B24-11-TN
% Mức: NB
Khi dòng điện chạy qua đoạn mạch ngoài nối giữa hai cực của nguồn điện thì các hạt mang điện trong mạch chuyển động có hướng dưới tác dụng của lực $\nguon{ly11 kntt.pdf}$
\choice
{Coulomb.}
{hấp dẫn.}
{lạ.}
{\True điện trường.}
\loigiai{
Chọn D.

Khi dòng điện chạy qua đoạn mạch ngoài nối giữa hai cực của nguồn điện thì các hạt mang điện trong mạch chuyển động có hướng dưới tác dụng của lực điện trường.
}
\end{ex}

% ===== Câu 19 =====
% ID: L11C4B24-12-TN
% Mức: NB
\begin{ex}
% ID: L11C4B24-12-TN
% Mức: NB
Khi dòng điện chạy qua nguồn điện thì các hạt mang điện ở bên trong nguồn điện chuyển động có hướng dưới tác dụng của lực $\nguon{ly11 kntt.pdf}$
\choice
{Coulomb.}
{hấp dẫn.}
{\True lạ.}
{điện trường.}
\loigiai{
Chọn C.

Khi dòng điện chạy qua nguồn điện thì các hạt mang điện ở bên trong nguồn điện chuyển động có hướng dưới tác dụng của lực lạ.
}
\end{ex}

% ===== Câu 20 =====
% ID: L11C4B24-13-TN
% Mức: NB
\begin{ex}
% ID: L11C4B24-13-TN
% Mức: NB
Một nguồn điện có suất điện động là $\mathcal{E},$ công của nguồn là $A$, độ lớn điện tích dịch chuyển qua nguồn là q. Mối liên hệ giữa các đại lượng này là $\nguon{ly11 kntt.pdf}$
\choice
{\True $A = q\mathcal{E}$.}
{$q = A\mathcal{E}$.}
{$\mathcal{E} = q$.}
{$A = q 2 \mathcal{E}$.}
\loigiai{
Chọn A.

Mối liên hệ giữa các đại lượng này là $A$ = qE
}
\end{ex}

% ===== Câu 21 =====
% ID: L11C4B24-14-TN
% Mức: NB
\begin{ex}
% ID: L11C4B24-14-TN
% Mức: NB
Công của lực lạ làm di chuyển điện tích $4\,\mathrm{C}$ từ cực âm đến cực dương bên trong nguồn điện là $24\,\mathrm{J}$. Suất điện động của nguồn là $\nguon{ly11 kntt.pdf}$
\choice
{$0,16\,\mathrm{V}$.}
{\True $6\,\mathrm{V}$.}
{$96\,\mathrm{V}$.}
{$0,6\,\mathrm{V}$.}
\loigiai{
Chọn B.

Suất điện động $\text{E} = \dfrac{A}{q} = \dfrac{24}{4} =$ 6 $\mathrm{V}$
}
\end{ex}

% ===== Câu 22 =====
% ID: L11C4B24-15-TN
% Mức: NB
\begin{ex}
% ID: L11C4B24-15-TN
% Mức: NB
Một nguồn điện có suất điện động $24\,\mathrm{V}$. Để chuyển một điện lượng $10\,\mathrm{C}$ qua nguồn thì lực lạ phải $\sin$ h một công là $\nguon{ly11 kntt.pdf}$
\choice
{$100\,\mathrm{J}$.}
{$2,4\,\mathrm{J}$.}
{$24\,\mathrm{J}$.}
{\True $240\,\mathrm{J}$.}
\loigiai{
Chọn D.

Công cần tìm $A= qE = 10.24 =240\,\mathrm{J}$
}
\end{ex}

% ===== Câu 23 =====
% ID: L11C4B24-16-TN
% Mức: NB
\begin{ex}
% ID: L11C4B24-16-TN
% Mức: NB
Một nguồn điện có suất điện động không đổi, để chuyển một điện lượng $5\,\mathrm{C}$ thì lực lạ phải $\sin$ h một công là 20 mJ. Để chuyển một điện lượng $10\,\mathrm{C}$ qua nguồn thì lực lạ phải $\sin$ h một công là $\nguon{ly11 kntt.pdf}$
\choice
{10 mJ.}
{15 mJ.}
{20 mJ.}
{\True 40 mJ.}
\loigiai{
Chọn D.

Công lực lạ $A = 10.\dfrac{20\cdot 10^{-3}}{5} = 40\cdot 10^{-3}J$
}
\end{ex}

% ===== Câu 24 =====
% ID: L11C4B24-17-TN
% Mức: NB
\begin{ex}
% ID: L11C4B24-17-TN
% Mức: NB
Một acquy có ghi thông số $12\,\mathrm{V}$ - 20Ah. Thông số này cho biết $\nguon{ly11 kntt.pdf}$
\choice
{điện lượng cực đại của acquy là $7200\,\mathrm{C}$.}
{điện trở trong của acquy là 0,16 $\Omega$.}
{\True dòng điện lớn nhất mà acquy có thể cung cấp là $20\,\mathrm{A}$.}
{năng lượng dự trữ của acquy là $12.106\,\mathrm{J}$.}
\loigiai{
Chọn C.

20 Ah có nghĩa là acquy cấp dòng điện $20\,\mathrm{A}$ trong vòng 1 giờ.

Điện lượng cực đại của acquy q = It = 20.3600 = $72000\,\mathrm{C}$.

Điện trở trong acquy $R = \dfrac{U}{I} = \dfrac{12}{20} = 0,6\Omega$

Năng lượng dự trữ $W$ = UIt = 12.20.3600 = $864000\,\mathrm{J}$.
}
\end{ex}

% ===== Câu 25 =====
% ID: L11C4B24-18-TN
% Mức: TH
\begin{ex}
% ID: L11C4B24-18-TN
% Mức: TH
Biểu thức tính công của nguồn điện có dòng điện không đổi là $\nguon{ly11 kntt.pdf}$
\choice
{$A = UIt$.}
{\True $A = Eit$}
{EIt - rI 2 t.}
{EIt + rI 2 t.}
\loigiai{
Chọn B.

Biểu thức công của nguồn điện có dòng điện không đổi là $A$ = EIt.
}
\end{ex}
